\documentclass[11pt,a4paper]{article}
\usepackage[utf8]{inputenc}
\usepackage[margin=1in]{geometry}
\usepackage{hyperref}
\usepackage{charter}
\usepackage{microtype}
\usepackage{booktabs}

\hypersetup{
    colorlinks=true,
    linkcolor=blue,
    urlcolor=blue,
    citecolor=blue
}

\setlength{\parindent}{0pt}
\setlength{\parskip}{6pt}

\begin{document}

\textbf{Seyed Mahmoud Sajjadi Mohammadabadi}\\
Department of Computer Science and Engineering\\
University of Nevada, Reno\\
1664 N. Virginia St., Reno, NV 89557, USA\\
Email: \href{mailto:mahmoud.sajjadi@unr.edu}{mahmoud.sajjadi@unr.edu}\\
ORCID: 0009-0001-9629-9734\\[1em]
\today

\textbf{Professor Gianfranco Chicco}\\
Editor-in-Chief, \textit{Sustainable Energy, Grids and Networks} (SEGAN)\\
Elsevier\\[1em]

\textbf{Subject: Submission of Original Research Article to \textit{Sustainable Energy, Grids and Networks}}\\[0.8em]

Dear Professor Chicco and Editorial Board,

I am pleased to submit my original research manuscript titled:

\begin{center}
\textbf{``Communication-Aware Event-Triggered Primary Frequency and Synthetic Inertia Support from Data Center UPS Fleets''}
\end{center}

for consideration as a Research Article in \textit{Sustainable Energy, Grids and Networks} (SEGAN).

\subsubsection*{Alignment with SEGAN Aims and Scope}
This manuscript directly advances several core focus areas specified in the SEGAN Aims and Scope, namely:
\begin{enumerate}
    \item \textbf{Control and planning of centralized and distributed energy storage}: Decentralized allocation of fast power reserves across multi-agent battery arrays while strictly respecting local state-of-charge boundaries and mission-critical backup runtime floors.
    \item \textbf{ICT integration in electrical and energy systems for monitoring and control}: Communication-aware event-triggering over local area networks with imperfect channels subject to packet drop, log-normal jitter, and variable sensing latencies.
    \item \textbf{Modeling of cyber-physical energy and communication systems}: Unified frequency-domain and time-domain stability formulations coupling physical power-frequency swing dynamics with discrete network protocols.
    \item \textbf{Demand-side management and demand response: definition of control schemes and protocols}: Harnessing pre-capitalized computing infrastructure to provide clean grid services without requiring dedicated utility-scale land footprint.
\end{enumerate}

\subsubsection*{Key Methodological and Practical Contributions}
The paper addresses the challenge of declining synchronous inertia in low-carbon grids by leveraging uninterruptible power supply (UPS) batteries in hyperscale data centers. Crucially, rather than presenting a case study or converter topology design, the paper provides fundamental methodological and computational contributions:
\begin{itemize}
    \item \textbf{Analytical Tracking and Nadir Degradation Bound}: We derive a closed-form analytical theorem (Section~2.3.1) proving that the worst-case frequency nadir degradation under adaptive event-triggering, network delay, and packet loss is strictly bounded by $|\Delta f_{\mathrm{err}}| \le 8.0\,\mathrm{mHz}$ relative to continuous droop.
    \item \textbf{Communication-Aware Adaptive Thresholding}: We formulate PULSE (Primary-frequency UPS Local Sub-second Event-triggered control), which eliminates idle-state communication under ambient drift and cuts post-contingency message traffic by $63.6\%$ compared to fixed-threshold schemes.
    \item \textbf{Analytical Delay Margin and Stability Proof}: Frequency-domain Nyquist and Bode stability analyses establish an open-loop phase margin of $68.2^\circ$ and an analytical delay margin $\tau_{\mathrm{margin}} = 1.58\,\mathrm{s}$, exceeding actual network delays by more than twenty-fold.
    \item \textbf{Sensing Realism and RoCoF Filtering}: We reframe sub-cycle synthetic inertia by leading with realistic 1-cycle PLL ($65.0\,\mathrm{ms}$) and 2-cycle PMU ($90.5\,\mathrm{ms}$) closed-loop power rise times, avoiding unphysical derivative chattering and inverter saturation.
    \item \textbf{Battery Autonomy and Degradation Assurance}: Dynamic battery modeling proves that providing 30~s grid support consumes only $0.166\,\mathrm{kWh}$ per rack, maintaining $>40\,\mathrm{min}$ of autonomous IT backup runtime ($100\%$ Tier III/IV compliance) with negligible 10-year cycling wear ($<0.018\%$).
\end{itemize}

\subsubsection*{Contribution to UN Sustainable Development Goals (SDGs)}
This work directly supports \textbf{UN SDG 7 (Affordable and Clean Energy)} and \textbf{SDG 11 (Sustainable Cities and Communities)} by transforming massive, pre-capitalized, and otherwise idle data center energy storage into low-carbon cyber-physical grid stability reserves, lowering reliance on fossil-fuel peaker plants.

\subsubsection*{Declarations and Submission Details}
\begin{itemize}
    \item \textbf{Sole Authorship}: The author is the sole contributor and conducted all conceptualization, modeling, simulation, formal analysis, and manuscript preparation.
    \item \textbf{Originality}: This manuscript is original, has not been published previously, and is not currently under consideration for publication elsewhere.
    \item \textbf{Conflict of Interest}: The author declares that there are no known competing financial interests or personal relationships that could have appeared to influence this work.
    \item \textbf{Data Availability}: All simulation code, automated verification suites, and numerical datasets are openly accessible via GitHub and permanently archived on Zenodo with DOI: \href{https://doi.org/10.5281/zenodo.14920188}{10.5281/zenodo.14920188}.
    \item \textbf{Generative AI Statement}: AI assisted with language editing and grammar polishing. The author reviewed and edited all content and takes full responsibility for the published article.
\end{itemize}

Thank you very much for your time and consideration of this manuscript.

\vspace{1em}
Sincerely,\\[1em]
\textbf{Seyed Mahmoud Sajjadi Mohammadabadi}\\
Department of Computer Science and Engineering\\
University of Nevada, Reno\\
Email: \href{mailto:mahmoudsajjadi@unr.edu}{mahmoudsajjadi@unr.edu}

\end{document}
